\documentclass[12pt,a4paper]{article}
\usepackage{fontspec}\setmainfont{Latin Modern Roman}\setsansfont{DejaVu Sans}
\usepackage{amsmath,amssymb,graphicx,geometry,enumitem,longtable,booktabs,fancyhdr,hyperref}
\geometry{left=28mm,right=23mm,top=24mm,bottom=24mm}
\hypersetup{unicode=true,colorlinks=true,linkcolor=black,urlcolor=blue,pdftitle={Giải tích Chương 1 Mục 1.3 Xây dựng hàm số mới từ những hàm số đã biết}}
\setlength{\parindent}{1.2em}\setlength{\parskip}{4pt}\setlength{\headheight}{16pt}
\pagestyle{fancy}\fancyhf{}\fancyhead[L]{\small Chương 1. Hàm số và mô hình}\fancyhead[R]{\small §1.3}\fancyfoot[C]{\thepage}
\newcommand{\fig}[3][.68]{\begin{center}\includegraphics[width=#1\linewidth]{assets/#2.png}\\{\small #3}\end{center}}
\newcommand{\ex}[1]{\par\medskip\noindent\textbf{Ví dụ #1.}\ }
\newcommand{\sol}{\par\noindent\textit{Lời giải.}\ }
\newcommand{\bt}[1]{\par\smallskip\noindent\textbf{#1.}\ }
\newcommand{\group}[1]{\par\medskip\noindent\textbf{#1}\par}
\begin{document}
\begin{center}
{\Large\bfseries GIẢI TÍCH}\\[5pt]
{\large Chương 1. Hàm số và mô hình}\\[9pt]
{\Large\bfseries 1.3. Xây dựng hàm số mới}\\[4pt]
{\large\bfseries từ những hàm số đã biết}
\end{center}
\noindent\textbf{Nguồn:} James Stewart, Daniel K. Clegg, Saleem Watson, \textit{Calculus: Early Transcendentals}, ấn bản thứ 9. Phạm vi: từ đầu §1.3 ở trang in 36 đến hết Bài tập 71 ở trang in 45 (trang PDF 71–80), trước §1.4.

\noindent\textbf{Quy ước biên tập:} Kế thừa thuật ngữ của §1.1–§1.2; giữ số ví dụ, hình và bài tập gốc. Công thức, bảng và các quy tắc được nhập lại bằng LaTeX. Hình đồ thị được trích từ sách; nhãn tiếng Anh còn trong hình có chú giải tiếng Việt. Số trang dịch được đánh riêng. Những lưu ý về cách diễn đạt hoặc sai sót của nguồn được tách riêng ở cuối mục.

Trong mục này, ta xuất phát từ các hàm cơ bản ở §1.2 để xây dựng những hàm số mới bằng cách tịnh tiến, co giãn và lấy đối xứng đồ thị. Ta cũng tìm hiểu cách kết hợp hai hàm số bằng các phép toán số học thông thường và phép hợp hàm.

\section*{Các phép biến đổi hàm số}
Áp dụng những phép biến đổi nhất định lên đồ thị của một hàm số đã cho, ta có thể thu được đồ thị của các hàm số liên quan. Điều này giúp phác họa nhanh bằng tay đồ thị của nhiều hàm số, đồng thời lập phương trình cho những đồ thị đã cho.

Trước hết, xét các phép tịnh tiến đồ thị. Nếu $c>0$, đồ thị $y=f(x)+c$ là đồ thị $y=f(x)$ được tịnh tiến lên trên $c$ đơn vị, vì mọi tung độ đều tăng cùng một lượng $c$. Tương tự, nếu $g(x)=f(x-c)$ với $c>0$, thì giá trị của $g$ tại $x$ bằng giá trị của $f$ tại $x-c$, tức tại một điểm nằm bên trái $x$ một khoảng $c$. Vì vậy, đồ thị $y=f(x-c)$ là đồ thị $y=f(x)$ được tịnh tiến sang phải $c$ đơn vị (Hình 1).

\group{Tịnh tiến theo phương thẳng đứng và phương ngang}
Giả sử $c>0$. Để có đồ thị ở cột trái, thực hiện phép biến đổi tương ứng đối với đồ thị $y=f(x)$:
\begin{center}\begin{tabular}{lp{.60\linewidth}}\toprule
Đồ thị cần dựng&Phép tịnh tiến\\\midrule
$y=f(x)+c$&Lên trên $c$ đơn vị.\\
$y=f(x)-c$&Xuống dưới $c$ đơn vị.\\
$y=f(x-c)$&Sang phải $c$ đơn vị.\\
$y=f(x+c)$&Sang trái $c$ đơn vị.\\\bottomrule
\end{tabular}\end{center}
\fig[.69]{f01}{Hình 1. Tịnh tiến đồ thị của $f$.}
Bây giờ xét các phép co giãn và lấy đối xứng. Nếu $c>1$, đồ thị $y=cf(x)$ thu được bằng cách giãn đồ thị $y=f(x)$ theo phương thẳng đứng với hệ số $c$, vì mỗi tung độ được nhân với $c$. Đồ thị $y=-f(x)$ là ảnh đối xứng của đồ thị $y=f(x)$ qua trục $x$, vì điểm $(x,y)$ được thay bằng $(x,-y)$. Hình 2 và bảng sau còn cho những phép co giãn, đối xứng khác.
\fig[.61]{f02}{Hình 2. Co giãn và lấy đối xứng đồ thị của $f$.}

\group{Co giãn theo phương thẳng đứng, phương ngang và lấy đối xứng}
Giả sử $c>1$. Thực hiện các biến đổi sau đối với đồ thị $y=f(x)$:
\begin{center}\begin{tabular}{lp{.60\linewidth}}\toprule
Đồ thị cần dựng&Phép biến đổi\\\midrule
$y=cf(x)$&Giãn theo phương thẳng đứng, tung độ gấp $c$ lần.\\
$y=\dfrac1c f(x)$&Co theo phương thẳng đứng, tung độ còn $1/c$ lần.\\[5pt]
$y=f(cx)$&Co theo phương ngang, hoành độ còn $1/c$ lần.\\
$y=f(x/c)$&Giãn theo phương ngang, hoành độ gấp $c$ lần.\\
$y=-f(x)$&Lấy đối xứng qua trục $x$.\\
$y=f(-x)$&Lấy đối xứng qua trục $y$.\\\bottomrule
\end{tabular}\end{center}
Hình 3 minh họa các phép co giãn trên đối với hàm côsin khi $c=2$. Chẳng hạn, muốn có đồ thị $y=2\cos x$, ta nhân tung độ mỗi điểm trên đồ thị $y=\cos x$ với 2; tức giãn đồ thị theo phương thẳng đứng với hệ số 2.
\fig[.93]{f03}{Hình 3. Các phép co giãn đồ thị hàm côsin.}

\ex{1} Từ đồ thị $y=\sqrt x$, dùng các phép biến đổi để vẽ đồ thị của
\[y=\sqrt x-2,\quad y=\sqrt{x-2},\quad y=-\sqrt x,\quad y=2\sqrt x,\quad y=\sqrt{-x}.\]
\sol Đồ thị hàm căn bậc hai $y=\sqrt x$, lấy từ Hình 1.2.13(a), được cho ở Hình 4(a). Các phần còn lại biểu diễn lần lượt: $y=\sqrt x-2$ bằng cách tịnh tiến xuống 2 đơn vị; $y=\sqrt{x-2}$ bằng cách tịnh tiến sang phải 2 đơn vị; $y=-\sqrt x$ bằng cách lấy đối xứng qua trục $x$; $y=2\sqrt x$ bằng cách giãn theo phương thẳng đứng với hệ số 2; $y=\sqrt{-x}$ bằng cách lấy đối xứng qua trục $y$.
\fig[1]{f04}{Hình 4. (a) $y=\sqrt x$; (b) $y=\sqrt x-2$; (c) $y=\sqrt{x-2}$;\\ (d) $y=-\sqrt x$; (e) $y=2\sqrt x$; (f) $y=\sqrt{-x}$.}

\ex{2} Phác họa đồ thị của $f(x)=x^2+6x+10$.
\sol Hoàn thành bình phương:
\[y=x^2+6x+10=(x+3)^2+1.\]
Vì vậy, xuất phát từ parabol $y=x^2$, ta tịnh tiến sang trái 3 đơn vị rồi lên trên 1 đơn vị để được đồ thị cần tìm (Hình 5).
\fig[.81]{f05}{Hình 5. (a) $y=x^2$; (b) $y=(x+3)^2+1$.}

\ex{3} Phác họa đồ thị của: (a) $y=\sin2x$; (b) $y=1-\sin x$.
\sol (a) Đồ thị $y=\sin2x$ thu được từ $y=\sin x$ bằng phép co theo phương ngang, hoành độ còn một nửa (Hình 6–7). Vì chu kỳ của $y=\sin x$ là $2\pi$, chu kỳ của $y=\sin2x$ là $2\pi/2=\pi$.
\fig[.74]{f06}{Hình 6. Đồ thị $y=\sin x$.}
\fig[.66]{f07}{Hình 7. Đồ thị $y=\sin2x$.}
(b) Xuất phát từ $y=\sin x$, lấy đối xứng qua trục $x$ để được $y=-\sin x$, rồi tịnh tiến lên trên 1 đơn vị để được $y=1-\sin x$ (Hình 8).
\fig[.78]{f08}{Hình 8. Đồ thị $y=1-\sin x$.}

\ex{4} Hình 9 cho đồ thị số giờ có ánh sáng ban ngày theo thời điểm trong năm tại một số vĩ độ. Biết Philadelphia ở khoảng vĩ độ $40^\circ$ Bắc, tìm hàm số mô hình hóa độ dài ban ngày tại thành phố này.
\fig[.75]{f09}{Hình 9. Độ dài ban ngày từ 21 tháng Ba đến 21 tháng Mười Hai ở các vĩ độ.\\ Hours: giờ; N: Bắc; Mar.–Dec.: các tháng Ba–Mười Hai.\\ Nguồn: phỏng theo L. Harrison, \textit{Daylight, Twilight, Darkness and Time}\\ (New York: Silver, Burdett, 1935), tr. 40.}
\sol Mỗi đường cong giống đồ thị hàm sin được tịnh tiến và co giãn. Từ đường màu xanh ứng với Philadelphia, độ dài ban ngày khoảng $14.8$ giờ vào ngày 21 tháng Sáu và $9.2$ giờ vào ngày 21 tháng Mười Hai. Vì vậy, biên độ, cũng là hệ số giãn đồ thị sin theo phương thẳng đứng, là
\[\frac12(14.8-9.2)=2.8.\]
Nếu thời gian $t$ được tính bằng ngày, cần giãn theo phương ngang với hệ số nào? Một năm có khoảng 365 ngày nên chu kỳ mô hình phải là 365, còn chu kỳ của $y=\sin t$ là $2\pi$. Bản gốc ghi hệ số giãn ngang ở đây là $2\pi/365$; xem ghi chú biên tập cuối mục để phân biệt hệ số trong đối số với hệ số giãn hình học.

Đường cong bắt đầu chu kỳ vào ngày 21 tháng Ba, ngày thứ 80 của năm, nên cần tịnh tiến sang phải 80 đơn vị; đồng thời tịnh tiến lên trên 12 đơn vị. Do đó, mô hình độ dài ban ngày tại Philadelphia vào ngày thứ $t$ trong năm là
\[L(t)=12+2.8\sin\left[\frac{2\pi}{365}(t-80)\right].\]

Một phép biến đổi đáng chú ý khác là lấy giá trị tuyệt đối của hàm số. Nếu $y=|f(x)|$, theo định nghĩa giá trị tuyệt đối, $y=f(x)$ khi $f(x)\ge0$, còn $y=-f(x)$ khi $f(x)<0$. Vì vậy, để thu được đồ thị $y=|f(x)|$ từ $y=f(x)$, ta giữ nguyên phần phía trên trục $x$ và lấy đối xứng phần phía dưới qua trục $x$.

\ex{5} Phác họa đồ thị $y=|x^2-1|$.
\sol Trước hết, vẽ parabol $y=x^2-1$ ở Hình 10(a) bằng cách tịnh tiến parabol $y=x^2$ xuống dưới 1 đơn vị. Đồ thị nằm dưới trục $x$ khi $-1<x<1$, nên ta lấy đối xứng phần đó qua trục $x$ để được đồ thị $y=|x^2-1|$ ở Hình 10(b).
\fig[.38]{f10}{Hình 10. (a) $y=x^2-1$; (b) $y=|x^2-1|$.}

\section*{Kết hợp các hàm số}
Hai hàm số $f$ và $g$ có thể kết hợp thành các hàm mới $f+g$, $f-g$, $fg$ và $f/g$, tương tự các phép cộng, trừ, nhân, chia số thực.
\begin{quote}\textbf{Định nghĩa.} Với hai hàm số $f$, $g$, các hàm tổng, hiệu, tích và thương được xác định bởi
\begin{align*}
(f+g)(x)&=f(x)+g(x),\\
(f-g)(x)&=f(x)-g(x),\\
(fg)(x)&=f(x)g(x),\\
\left(\frac fg\right)(x)&=\frac{f(x)}{g(x)}.
\end{align*}\end{quote}
Nếu tập xác định của $f$ là $A$, của $g$ là $B$, thì tập xác định của $f+g$ và $f-g$ là giao $A\cap B$, vì cả $f(x)$ lẫn $g(x)$ đều phải xác định. Chẳng hạn, $f(x)=\sqrt x$ có tập xác định $A=[0,\infty)$, $g(x)=\sqrt{2-x}$ có tập xác định $B=(-\infty,2]$; do đó,
\[(f+g)(x)=\sqrt x+\sqrt{2-x}\]
có tập xác định $A\cap B=[0,2]$.

Tập xác định của $fg$ cũng là $A\cap B$. Vì không được chia cho 0, tập xác định của $f/g$ là
\[\{x\in A\cap B\mid g(x)\ne0\}.\]
Ví dụ, nếu $f(x)=x^2$ và $g(x)=x-1$, hàm hữu tỉ $(f/g)(x)=x^2/(x-1)$ có tập xác định $\{x\mid x\ne1\}$, hay $(-\infty,1)\cup(1,\infty)$.

Có một cách kết hợp hai hàm số khác. Giả sử $y=f(u)=\sqrt u$ và $u=g(x)=x^2+1$. Vì $y$ là hàm của $u$, còn $u$ là hàm của $x$, nên cuối cùng $y$ là hàm của $x$. Bằng phép thế,
\[y=f(u)=f(g(x))=f(x^2+1)=\sqrt{x^2+1}.\]
Quy trình này gọi là \textbf{phép hợp hàm}, vì hàm mới được xây dựng bằng cách hợp hai hàm đã cho $f$ và $g$.

Nói chung, với hai hàm $f$, $g$ bất kỳ, ta lấy $x$ trong tập xác định của $g$ và tính $g(x)$. Nếu $g(x)$ thuộc tập xác định của $f$, ta tính tiếp $f(g(x))$. Đầu ra của hàm trước được dùng làm đầu vào của hàm sau. Hàm số mới $h(x)=f(g(x))$ thu được bằng cách thế $g$ vào $f$ được gọi là \textbf{hàm hợp} của $f$ và $g$, ký hiệu $f\circ g$, đọc là ``$f$ hợp $g$''.
\begin{quote}\textbf{Định nghĩa.} Với hai hàm số $f$, $g$, hàm hợp $f\circ g$ được xác định bởi
\[(f\circ g)(x)=f(g(x)).\]
\end{quote}
Tập xác định của $f\circ g$ gồm mọi $x$ thuộc tập xác định của $g$ sao cho $g(x)$ thuộc tập xác định của $f$. Nói cách khác, $(f\circ g)(x)$ xác định khi cả $g(x)$ và $f(g(x))$ đều xác định. Hình 11 minh họa hàm hợp bằng sơ đồ máy.
\fig[.36]{f11}{Hình 11. Máy $f\circ g$ gồm máy $g$ thực hiện trước, rồi máy $f$.\\ input: đầu vào; output: đầu ra.}

\ex{6} Cho $f(x)=x^2$, $g(x)=x-3$. Tìm $f\circ g$ và $g\circ f$.
\sol Ta có
\begin{align*}
(f\circ g)(x)&=f(g(x))=f(x-3)=(x-3)^2,\\
(g\circ f)(x)&=g(f(x))=g(x^2)=x^2-3.
\end{align*}
\noindent\textbf{Lưu ý.} Ví dụ 6 cho thấy nói chung $f\circ g\ne g\circ f$. Ký hiệu $f\circ g$ nghĩa là thực hiện $g$ trước, rồi $f$. Trong ví dụ, $f\circ g$ là trừ 3 rồi bình phương; $g\circ f$ là bình phương rồi trừ 3.

\ex{7} Cho $f(x)=\sqrt x$, $g(x)=\sqrt{2-x}$. Tìm công thức và tập xác định của: (a) $f\circ g$; (b) $g\circ f$; (c) $f\circ f$; (d) $g\circ g$.
\sol (a)
\[(f\circ g)(x)=f(g(x))=f(\sqrt{2-x})=\sqrt{\sqrt{2-x}}=\sqrt[4]{2-x}.\]
Tập xác định là $\{x\mid2-x\ge0\}=\{x\mid x\le2\}=(-\infty,2]$.

(b)
\[(g\circ f)(x)=g(f(x))=g(\sqrt x)=\sqrt{2-\sqrt x}.\]
Để $\sqrt x$ xác định, cần $x\ge0$. Để $\sqrt{2-\sqrt x}$ xác định, cần $2-\sqrt x\ge0$, tức $\sqrt x\le2$, hay $x\le4$. Vậy $0\le x\le4$, tập xác định là $[0,4]$.

\noindent\textit{Chú thích bên lề của bản gốc:} Nếu $0\le a\le b$ thì $a^2\le b^2$.

(c)
\[(f\circ f)(x)=f(f(x))=f(\sqrt x)=\sqrt{\sqrt x}=\sqrt[4]x.\]
Tập xác định là $[0,\infty)$.

(d)
\[(g\circ g)(x)=g(g(x))=g(\sqrt{2-x})=\sqrt{2-\sqrt{2-x}}.\]
Biểu thức xác định khi đồng thời $2-x\ge0$ và $2-\sqrt{2-x}\ge0$. Điều kiện thứ nhất là $x\le2$; điều kiện thứ hai tương đương $\sqrt{2-x}\le2$, hay $2-x\le4$, tức $x\ge-2$. Do đó $-2\le x\le2$, tập xác định là $[-2,2]$.

Có thể hợp ba hàm số trở lên. Chẳng hạn, $f\circ g\circ h$ được tính bằng cách thực hiện $h$ trước, rồi $g$, cuối cùng $f$:
\[(f\circ g\circ h)(x)=f(g(h(x))).\]

\ex{8} Tìm $f\circ g\circ h$ nếu $f(x)=x/(x+1)$, $g(x)=x^{10}$ và $h(x)=x+3$.
\sol
\begin{align*}
(f\circ g\circ h)(x)&=f(g(h(x)))=f(g(x+3))\\
&=f((x+3)^{10})=\frac{(x+3)^{10}}{(x+3)^{10}+1}.
\end{align*}
Đến đây, ta đã dùng phép hợp để xây dựng hàm phức tạp từ những hàm đơn giản. Trong giải tích, việc phân tích một hàm phức tạp thành các hàm đơn giản cũng thường hữu ích, như ví dụ tiếp theo.

\ex{9} Cho $F(x)=\cos^2(x+9)$. Tìm $f$, $g$, $h$ sao cho $F=f\circ g\circ h$.
\sol Vì $F(x)=[\cos(x+9)]^2$, công thức cho biết: trước hết cộng 9, sau đó lấy côsin, cuối cùng bình phương. Vì vậy, đặt
\[h(x)=x+9,\qquad g(x)=\cos x,\qquad f(x)=x^2.\]
Khi đó
\begin{align*}
(f\circ g\circ h)(x)&=f(g(h(x)))=f(g(x+9))\\
&=f(\cos(x+9))=[\cos(x+9)]^2=F(x).
\end{align*}

\clearpage
\section*{1.3. Bài tập}
\bt{1} Giả sử đã cho đồ thị $f$. Viết phương trình của đồ thị nhận được khi: (a) Tịnh tiến lên 3 đơn vị. (b) Tịnh tiến xuống 3 đơn vị. (c) Tịnh tiến sang phải 3 đơn vị. (d) Tịnh tiến sang trái 3 đơn vị. (e) Lấy đối xứng qua trục $x$. (f) Lấy đối xứng qua trục $y$. (g) Giãn theo phương thẳng đứng với hệ số 3. (h) Co theo phương thẳng đứng, tung độ còn $1/3$ lần.
\bt{2} Giải thích cách thu được mỗi đồ thị sau từ $y=f(x)$: (a) $y=f(x)+8$; (b) $y=f(x+8)$; (c) $y=8f(x)$; (d) $y=f(8x)$; (e) $y=-f(x)-1$; (f) $y=8f(x/8)$.
\bt{3} Cho đồ thị $y=f(x)$. Ghép mỗi phương trình với đồ thị tương ứng và giải thích: (a) $y=f(x-4)$; (b) $y=f(x)+3$; (c) $y=\frac13f(x)$; (d) $y=-f(x+4)$; (e) $y=2f(x+6)$.
\fig[.76]{e03}{Đồ thị cho Bài tập 3, với các nhãn số 1–5 khoanh tròn theo bản gốc.}
\bt{4} Từ đồ thị $f$ đã cho, vẽ: (a) $y=f(x)-3$; (b) $y=f(x+1)$; (c) $y=\frac12f(x)$; (d) $y=-f(x)$.
\fig[.40]{e04}{Đồ thị cho Bài tập 4.}
\bt{5} Dùng đồ thị $f$ đã cho để vẽ: (a) $y=f(2x)$; (b) $y=f(x/2)$; (c) $y=f(-x)$; (d) $y=-f(-x)$.
\fig[.56]{e05}{Đồ thị cho Bài tập 5.}
\group{6–7. Cho đồ thị $y=\sqrt{3x-x^2}$. Dùng các phép biến đổi để tạo hàm số có đồ thị như các hình yêu cầu.}
\fig[.56]{e06_07_base}{Đồ thị xuất phát cho Bài tập 6–7.}
\fig[.86]{e06_07}{Các đồ thị cần tạo trong Bài tập 6 (trái) và 7 (phải).}
\bt{8} (a) Đồ thị $y=1+\sqrt x$ liên hệ thế nào với đồ thị $y=\sqrt x$? Dùng câu trả lời và Hình 4(a) để phác họa $y=1+\sqrt x$. (b) Đồ thị $y=5\sin\pi x$ liên hệ thế nào với $y=\sin x$? Dùng câu trả lời và Hình 6 để phác họa $y=5\sin\pi x$.
\group{9–26. Vẽ đồ thị bằng tay bằng cách xuất phát từ một hàm cơ bản ở Bảng 1.2.3, rồi áp dụng các phép biến đổi thích hợp; không vẽ bằng cách lập bảng điểm.}
\bt{9} $y=1+x^2$.
\bt{10} $y=(x+1)^2$.
\bt{11} $y=|x+2|$.
\bt{12} $y=1-x^3$.
\bt{13} $y=\dfrac1x+2$.
\bt{14} $y=-\sqrt x-1$.
\bt{15} $y=\sin4x$.
\bt{16} $y=1+\dfrac1{x^2}$.
\bt{17} $y=2+\sqrt{x+1}$.
\bt{18} $y=-(x-1)^2+3$.
\bt{19} $y=x^2-2x+5$.
\bt{20} $y=(x+1)^3+2$.
\bt{21} $y=2-|x|$.
\bt{22} $y=2-2\cos x$.
\bt{23} $y=3\sin\frac12x+1$.
\bt{24} $y=\frac14\tan\left(x-\frac\pi4\right)$.
\bt{25} $y=|\cos\pi x|$.
\bt{26} $y=|\sqrt x-1|$.
\bt{27} New Orleans ở vĩ độ $30^\circ$ Bắc. Dùng Hình 9 để tìm hàm số mô hình hóa số giờ có ánh sáng ban ngày tại đây theo thời điểm trong năm. Kiểm tra độ chính xác bằng thông tin: ngày 31 tháng Ba, Mặt Trời mọc lúc 5 giờ 51 phút và lặn lúc 18 giờ 18 phút.
\bt{28} Sao biến quang là ngôi sao có độ sáng tăng giảm luân phiên. Với Delta Cephei, ngôi sao biến quang dễ quan sát nhất, thời gian giữa hai lần sáng cực đại là $5.4$ ngày; độ sáng trung bình (cấp sao) là $4.0$, biến thiên $\pm0.35$ cấp sao. Tìm hàm mô hình hóa độ sáng của Delta Cephei theo thời gian.
\bt{29} Một số mức thủy triều cao nhất thế giới xuất hiện tại vịnh Fundy, bờ Đại Tây Dương của Canada. Ở Hopewell Cape, độ sâu nước khi triều thấp khoảng $2.0$ m, khi triều cao khoảng $12.0$ m. Chu kỳ dao động tự nhiên khoảng 12 giờ; trong một ngày cụ thể, triều cao xảy ra lúc 6 giờ 45 phút. Tìm hàm có chứa côsin mô hình hóa độ sâu $D(t)$ (m) theo thời gian $t$ (giờ kể từ nửa đêm) trong ngày đó.
\bt{30} Trong một chu kỳ hô hấp bình thường, thể tích khí ra vào phổi khoảng 500 mL. Thể tích khí dự trữ và khí cặn còn trong phổi khoảng 2000 mL; một chu kỳ hô hấp ở người trung bình kéo dài khoảng 4 giây. Tìm mô hình cho tổng thể tích khí $V(t)$ trong phổi theo thời gian.
\bt{31} (a) Đồ thị $y=f(|x|)$ liên hệ thế nào với đồ thị $f$? (b) Phác họa $y=\sin|x|$. (c) Phác họa $y=\sqrt{|x|}$.
\bt{32} Dùng đồ thị $f$ đã cho để phác họa $y=1/f(x)$. Những đặc điểm nào của $f$ quan trọng nhất đối với việc phác họa? Giải thích cách sử dụng chúng.
\fig[.48]{e32}{Đồ thị cho Bài tập 32.}
\group{33–34. Tìm (a) $f+g$; (b) $f-g$; (c) $fg$; (d) $f/g$ và nêu tập xác định.}
\bt{33} $f(x)=\sqrt{25-x^2}$, $g(x)=\sqrt{x+1}$.
\bt{34} $f(x)=\dfrac1{x-1}$, $g(x)=\dfrac1x-2$.
\group{35–40. Tìm công thức và tập xác định của (a) $f\circ g$; (b) $g\circ f$; (c) $f\circ f$; (d) $g\circ g$.}
\bt{35} $f(x)=x^3+5$, $g(x)=\sqrt[3]x$.
\bt{36} $f(x)=\dfrac1x$, $g(x)=2x+1$.
\bt{37} $f(x)=\dfrac1{\sqrt x}$, $g(x)=x+1$.
\bt{38} $f(x)=\dfrac{x}{x+1}$, $g(x)=2x-1$.
\bt{39} $f(x)=\dfrac2x$, $g(x)=\sin x$.
\bt{40} $f(x)=\sqrt{5-x}$, $g(x)=\sqrt{x-1}$.
\group{41–44. Tìm $f\circ g\circ h$.}
\bt{41} $f(x)=3x-2$, $g(x)=\sin x$, $h(x)=x^2$.
\bt{42} $f(x)=|x-4|$, $g(x)=2^x$, $h(x)=\sqrt x$.
\bt{43} $f(x)=\sqrt{x-3}$, $g(x)=x^2$, $h(x)=x^3+2$.
\bt{44} $f(x)=\tan x$, $g(x)=\dfrac{x}{x-1}$, $h(x)=\sqrt[3]x$.
\group{45–50. Biểu diễn hàm số dưới dạng $f\circ g$.}
\bt{45} $F(x)=(2x+x^2)^4$.
\bt{46} $F(x)=\cos^2x$.
\bt{47} $F(x)=\dfrac{\sqrt[3]x}{1+\sqrt[3]x}$.
\bt{48} $G(x)=\sqrt[3]{\dfrac{x}{1+x}}$.
\bt{49} $v(t)=\sec(t^2)\tan(t^2)$.
\bt{50} $H(x)=\sqrt{1+\sqrt x}$.
\group{51–54. Biểu diễn hàm số dưới dạng $f\circ g\circ h$.}
\bt{51} $R(x)=\sqrt{\sqrt x-1}$.
\bt{52} $H(x)=\sqrt[8]{2+|x|}$.
\bt{53} $S(t)=\sin^2(\cos t)$.
\bt{54} $H(t)=\cos\left(\sqrt{\tan t}+1\right)$.
\group{55–56. Dùng bảng sau để tính các biểu thức.}
\begin{center}\begin{tabular}{c|rrrrrr}\toprule
$x$&1&2&3&4&5&6\\\midrule
$f(x)$&3&1&5&6&2&4\\
$g(x)$&5&3&4&1&3&2\\\bottomrule\end{tabular}\end{center}
\bt{55} (a) $f(g(3))$; (b) $g(f(2))$; (c) $(f\circ g)(5)$; (d) $(g\circ f)(5)$.
\bt{56} (a) $g(g(g(2)))$; (b) $(f\circ f\circ f)(1)$; (c) $(f\circ f\circ g)(1)$; (d) $(g\circ f\circ g)(3)$.
\bt{57} Dùng đồ thị $f$ và $g$ để tính các biểu thức sau, hoặc giải thích vì sao không xác định: (a) $f(g(2))$; (b) $g(f(0))$; (c) $(f\circ g)(0)$; (d) $(g\circ f)(6)$; (e) $(g\circ g)(-2)$; (f) $(f\circ f)(4)$.
\fig[.56]{e57}{Đồ thị cho Bài tập 57.}
\bt{58} Dùng đồ thị $f$ và $g$ để ước lượng $f(g(x))$ với $x=-5,-4,-3,\ldots,5$. Từ các ước lượng đó, phác họa đồ thị $f\circ g$.
\fig[.56]{e58}{Đồ thị cho Bài tập 58.}
\bt{59} Một hòn đá rơi xuống hồ tạo gợn sóng tròn lan ra với tốc độ 60 cm/s. (a) Biểu diễn bán kính $r$ của vòng tròn theo thời gian $t$ (s). (b) Nếu $A$ là diện tích hình tròn theo bán kính, tìm $A\circ r$ và diễn giải ý nghĩa.
\bt{60} Một quả bóng hình cầu được bơm phồng, bán kính tăng với tốc độ 2 cm/s. (a) Biểu diễn bán kính $r$ theo thời gian $t$ (s). (b) Nếu $V$ là thể tích theo bán kính, tìm $V\circ r$ và diễn giải ý nghĩa.
\bt{61} Một con tàu chạy với tốc độ 30 km/h song song với bờ biển thẳng, cách bờ 6 km và đi ngang một ngọn hải đăng lúc giữa trưa. (a) Biểu diễn khoảng cách $s$ từ hải đăng đến tàu theo quãng đường $d$ tàu đã đi kể từ giữa trưa; tức tìm $f$ sao cho $s=f(d)$. (b) Biểu diễn $d$ theo thời gian $t$ kể từ giữa trưa; tức tìm $g$ sao cho $d=g(t)$. (c) Tìm $f\circ g$. Hàm số này biểu thị điều gì?
\bt{62} Một máy bay bay với tốc độ 350 dặm/giờ, ở độ cao 1 dặm, đi qua đúng phía trên trạm radar lúc $t=0$. (a) Biểu diễn quãng đường theo phương ngang $d$ (dặm) mà máy bay đã bay theo $t$. (b) Biểu diễn khoảng cách $s$ từ máy bay đến trạm radar theo $d$. (c) Dùng phép hợp để biểu diễn $s$ theo $t$.
\bt{63} \textbf{Hàm Heaviside.} Hàm $H$ được xác định bởi
\[H(t)=\begin{cases}0&\text{nếu }t<0,\\1&\text{nếu }t\ge0.\end{cases}\]
Hàm được dùng khi nghiên cứu mạch điện để biểu diễn sự tăng đột ngột của dòng điện hoặc điện áp khi đóng công tắc tức thời. (a) Phác họa đồ thị $H$. (b) Phác họa điện áp $V(t)$ nếu công tắc đóng lúc $t=0$ và điện áp 120 V được đặt tức thời vào mạch. Viết $V(t)$ theo $H(t)$. (c) Phác họa $V(t)$ nếu công tắc đóng lúc $t=5$ s và điện áp 240 V được đặt tức thời vào mạch. Viết công thức theo $H(t)$; lưu ý thời điểm bắt đầu $t=5$ tương ứng với một phép tịnh tiến.
\bt{64} \textbf{Hàm dốc.} Hàm Heaviside ở Bài tập 63 có thể dùng để xác định hàm dốc $y=ctH(t)$, biểu diễn sự tăng dần của điện áp hoặc dòng điện. (a) Phác họa $y=tH(t)$. (b) Phác họa $V(t)$ nếu công tắc đóng lúc $t=0$ và điện áp tăng dần đến 120 V trong 60 s. Viết công thức theo $H(t)$ với $t\le60$. (c) Phác họa $V(t)$ nếu công tắc đóng lúc $t=7$ s và điện áp tăng dần đến 100 V trong 25 s. Viết công thức theo $H(t)$ với $t\le32$.
\bt{65} Cho hai hàm tuyến tính $f(x)=m_1x+b_1$ và $g(x)=m_2x+b_2$. Hàm $f\circ g$ có tuyến tính không? Nếu có, hệ số góc của đồ thị là bao nhiêu?
\bt{66} Đầu tư $x$ đô la với lãi suất 4\%, ghép lãi hằng năm, thì sau một năm có $A(x)=1.04x$. Tìm $A\circ A$, $A\circ A\circ A$ và $A\circ A\circ A\circ A$. Các hàm hợp biểu thị điều gì? Tìm công thức khi hợp $n$ bản sao của $A$.
\bt{67} (a) Cho $g(x)=2x+1$, $h(x)=4x^2+4x+7$. Tìm $f$ sao cho $f\circ g=h$. Hãy nghĩ về các phép toán cần thực hiện trên công thức của $g$ để thu được công thức của $h$. (b) Cho $f(x)=3x+5$, $h(x)=3x^2+3x+2$. Tìm $g$ sao cho $f\circ g=h$.
\bt{68} Cho $f(x)=x+4$ và $h(x)=4x-1$. Tìm $g$ sao cho $g\circ f=h$.
\bt{69} Giả sử $g$ là hàm số chẵn và $h=f\circ g$. Hàm $h$ có luôn chẵn không?
\bt{70} Giả sử $g$ là hàm số lẻ và $h=f\circ g$. Hàm $h$ có luôn lẻ không? Nếu $f$ lẻ thì sao? Nếu $f$ chẵn thì sao?
\bt{71} Cho hàm số $f(x)$ có tập xác định $\mathbb R$. (a) Chứng minh $E(x)=f(x)+f(-x)$ chẵn. (b) Chứng minh $O(x)=f(x)-f(-x)$ lẻ. (c) Chứng minh mọi hàm $f(x)$ đều viết được thành tổng một hàm chẵn và một hàm lẻ. (d) Biểu diễn $f(x)=2^x+(x-3)^2$ thành tổng một hàm chẵn và một hàm lẻ.

\clearpage
\section*{Ghi chú biên tập}
\noindent\textit{Nội dung bổ sung, không thuộc nguyên bản.}
\begin{enumerate}[leftmargin=*,itemsep=5pt]
\item Trong Ví dụ 4, bản gốc ghi hệ số giãn ngang là $2\pi/365$. Cần phân biệt: \textbf{hệ số nhân biến trong đối số sin} là $2\pi/365$, còn \textbf{hệ số giãn hình học theo phương ngang} là nghịch đảo của nó, $365/(2\pi)$. Công thức mô hình $L(t)$ của sách là đúng và được giữ nguyên.
\item Phép co theo hệ số $c>1$ được diễn đạt là kích thước tương ứng còn $1/c$ lần, để tránh nhầm với phép giãn $c$ lần.
\item Phân biệt $(fg)(x)=f(x)g(x)$ với $(f\circ g)(x)=f(g(x))$. Với hàm hợp, thực hiện hàm ở bên phải trước; tập xác định phải được kiểm tra trước khi rút gọn công thức.
\item Bài tập 63 dùng quy ước $H(0)=1$. Trong các tài liệu khác, giá trị tại 0 có thể được quy ước khác; bản dịch giữ đúng định nghĩa của đề gốc.
\end{enumerate}
\section*{Thuật ngữ bổ sung từ §1.3}
\begingroup\small
\begin{longtable}{p{.43\linewidth}p{.49\linewidth}}\toprule Tiếng Anh&Tiếng Việt thống nhất\\\midrule\endhead
transformation of a graph&phép biến đổi đồ thị\\
translation / shift&phép tịnh tiến\\
vertical shift&tịnh tiến theo phương thẳng đứng\\
horizontal shift&tịnh tiến theo phương ngang\\
stretching&phép giãn\\
shrinking / compression&phép co\\
vertical stretching&giãn theo phương thẳng đứng\\
horizontal stretching&giãn theo phương ngang\\
reflection&phép lấy đối xứng\\
reflection about the x-axis&đối xứng qua trục hoành\\
reflection about the y-axis&đối xứng qua trục tung\\
completing the square&hoàn thành bình phương\\
amplitude&biên độ\\
sum of functions&hàm tổng\\
difference of functions&hàm hiệu\\
product of functions&hàm tích\\
quotient of functions&hàm thương\\
intersection&giao\\
composition&phép hợp hàm\\
composite function&hàm hợp\\
decomposition of a function&phân tích hàm số thành các hàm thành phần\\
Heaviside function&hàm Heaviside\\
ramp function&hàm dốc\\
\bottomrule\end{longtable}
\endgroup
\noindent\textbf{Phạm vi hoàn thành:} Trọn §1.3, gồm Ví dụ 1–9, Hình 1–11 và Bài tập 1–71. Hai bảng quy tắc biến đổi và bảng giá trị cho Bài tập 55–56 được nhập lại đầy đủ. Phần tiếp theo: §1.4, \textit{Hàm số mũ}.
\end{document}
